\section{In the Midst of Silence: The Sunday within the Octave}\label{sec:sunday-octave}

Whenever 25~December falls on Monday through Saturday---eighty-six years in
every century---a Sunday falls inside the Christmas Octave between 26 and
31~December. Under the 1960 \textit{Codex Rubricarum} (General Rubric~69 and
RG~91, row~15), that Sunday is always celebrated as the \textbf{Sunday within
the Octave of the Nativity}, a second-class Sunday in white vestments with
\textit{Gloria}, Creed, the Nativity Preface, and the proper
\textit{Communicantes} (\textit{diem sacratissimum}).

Its precedence reshapes the Octave differently depending on the date it hits.
When Sunday falls on \textbf{26, 27, or 28~December}, it outranks St.~Stephen,
St.~John, or the Holy Innocents (row~15 over row~16), admitting the saint's
feast as a single ordinary commemoration in Low Mass only (omitted in a
non-conventual Sung Mass, RG~111~b). No separate commemoration of the Octave is
added, because the Sunday itself belongs to the Octave. When Sunday falls on
\textbf{29, 30, or 31~December}, it replaces the fifth, sixth, or seventh
Octave day outright, and the ordinary commemorations of St.~Thomas of Canterbury
(29~December) and St.~Sylvester~I (31~December) are omitted entirely, even at
Low Mass. Finally, in the fifteen years per century when \textbf{25~December is
itself a Sunday} (such as 2022, 2033, and 2039), 1~January is also a Sunday; no
Sunday falls between them, and the Mass \textit{Dum medium silentium} is not
celebrated anywhere in that liturgical year.

In the Divine Office, this Sunday also illustrates one of the 1960 Breviary
reforms: under Breviary Rubrics~162~a and~220~b, Matins has one nocturn of nine
psalms and three lessons---two from the opening of St.~Paul's Epistle to the
Romans (Rom~1) and a third taken solely from the opening of St.~Ambrose's Gospel
homily on Simeon's prophecy (Luke~2:33, the former seventh lesson)---whereas
older pre-1960 Breviaries sang three nocturns and included St.~Leo the Great's
ninth Nativity sermon together with St.~Ambrose's ninth lesson on the prophetess
Anna.

\subsection{The Threefold Silence Broken by the Word}

Few Introits in the entire Roman Missal arrest the ear like Wisdom~18:14--15:
\textit{Dum medium silentium tenerent omnia, et nox in suo cursu medium iter
haberet, omnipotens Sermo tuus, Domine, de caelis a regalibus sedibus venit},
``While all things were in quiet silence, and the night was in the midst of her
course, thy Almighty word, O Lord, came down from heaven from thy royal
throne.'' In the Book of Wisdom, those lines describe the Angel of the Lord
leaping down into Egypt on the Passover night; the Church takes them onto her
lips for the greater Passover of the Incarnation, joining them to
Psalm~92~[93]:1, ``The Lord hath reigned, he is clothed with beauty.''

In Book~VI of his \textit{Rationale}, William Durandus explains why the Church
sings of ``mid-silence'' on this Sunday: there is a threefold silence,
\textit{scilicet ignorantiae, desperationis, et gloriae}---the silence of human
ignorance when the world knew not God, the silence of despair under sin's
bondage, and the silence of contemplative wonder in glory. Into the midnight of
a silent world, the Word leaps from the Father's throne to speak life. The
Collect then turns that descent into daily faithfulness:

\begin{massbox}{Collect of the Sunday within the Octave of the Nativity}
\textbf{Latin} (1962 Roman Missal). Omnipotens sempiterne Deus, dirige actus
nostros in beneplacito tuo: ut in nomine dilecti Filii tui mereamur bonis
operibus abundare: Qui tecum vivit et regnat in unitate.

\textbf{English} (Lasance, 1945). O almighty and eternal God, direct our actions
in conformity with Thy good pleasure, that in the name of Thy beloved Son we may
be worthy to abound in good works. Who with Thee.
\end{massbox}

\subsection{Adoption as Sons and the Sword Foretold}

The Epistle (Gal~4:1--7) unfolds the fruit of the Incarnation in six crisp
verses: ``But when the fulness of the time was come, God sent his Son, made of a
woman, made under the law: That he might redeem them who were under the law:
that we might receive the adoption of sons. And because you are sons, God hath
sent the Spirit of his Son into your hearts, crying: Abba, Father. Therefore,
now he is not a servant, but a son. And if a son, an heir also through God.''
The Gradual answers with Psalm~44~[45]:3, \textit{Speciosus forma prae filiis
hominum}, ``Thou art beautiful above the sons of men: grace is poured abroad in
thy lips''---a bridal psalm of the King that the Offertory (\textit{Deus
firmavit orbem terrae}, Ps~92~[93]:1--2) anchors in the unshakable foundation
of his throne.

Then the Gospel (Luke~2:33--40) carries the assembly from the stable to the
Temple in Jerusalem, where Mary and Joseph stand wondering at Simeon's words:

\begin{quote}
And Simeon blessed them and said to Mary his mother: Behold this child is set
for the fall and for the resurrection of many in Israel and for a sign which
shall be contradicted. And thy own soul a sword shall pierce, that, out of many
hearts thoughts may be revealed.
\end{quote}

Even inside the bright white vestments of the Christmas Octave, the shadow of
Good Friday falls squarely across Mary's heart. The Child who brings peace to
men of good will is also ``a sign which shall be contradicted,'' dividing every
heart that meets him.

In the third lesson of 1962 Sunday Matins (from Book~II of his
\textit{Exposition of Luke}), St.~Ambrose comments directly on Simeon's words:
``We see that God's abounding grace is poured forth on all by the birth of the
Lord, and that the gift of prophecy is not denied to the righteous, but to the
unbelieving. Simeon prophesieth that our Lord Jesus Christ is set for the fall
and rising again of many in Israel, setting forth that the just and the unjust
reap different fruits from the coming of the Saviour.''

Immediately beside the aged Simeon stands a figure too often passed over in
haste: ``one Anna, a prophetess, the daughter of Phanuel, of the tribe of Aser,''
a widow of eighty-four years ``who departed not from the temple, by fastings and
prayers serving night and day,'' who comes in at that very hour to confess the
Lord and speak of him to all who looked for the redemption of Israel
(Luke~2:36--38). Earlier in the same commentary (Book~II, section~58), Ambrose
marvels that every age and both sexes build up the Church's faith---a virgin
conceiving, a barren woman bearing, a mute man speaking, Elizabeth prophesying,
a Magus adoring, an unborn child leaping in the womb, a widow confessing, and a
just man awaiting the Lord. And in section~62 (formerly the ninth lesson in
pre-1960 Breviaries), he crowns Anna's witness: ``There had been a triple
prophecy; the prophecy of Simeon had followed the prophecy of the virgin, and
the prophecy of the wife; those, namely, of Mary and Elizabeth. And now ought
the widow also to prophesy, that no sex nor state might be wanting''
(\textit{ne qua aut professio deesset, aut sexus}).

At Communion, the choir sings the angel's words to Joseph after Herod's death:
\textit{Tolle Puerum et Matrem ejus, et vade in terram Israel}, ``Arise, and
take the child and his mother, and go into the land of Israel. For they are dead
that sought the life of the child'' (Matt~2:20)---binding the Sunday straight
back to the Flight into Egypt and the Holy Innocents. And the Gospel ends in
quiet domestic growth at Nazareth: ``And the child grew and waxed strong, full
of wisdom: and the grace of God was in him'' (Luke~2:40). In the middle of
Christmas week, the Sunday within the Octave asks the Christian household to
carry the silence of Bethlehem back into ordinary days, letting the grace of the
Son's adoption direct every act ``in conformity with Thy good pleasure.''
